% ── Chapter 3: Methodology ────────────────────────────────────
\chapter{ระเบียบวิธีวิจัย (Methodology)}

\section{ภาพรวมสถาปัตยกรรมระบบ (System Architecture Overview)}

ระบบที่นำเสนอคือ Full-stack web application ที่ใช้สถาปัตยกรรมแบบ Client-server ซึ่งออกแบบตามหลักการแยกส่วนความรับผิดชอบ (Separation of concerns) ระหว่างส่วนติดต่อผู้ใช้ (User interface), AI inference pipeline และ Data persistence layer การแยกส่วนทางสถาปัตยกรรมนี้ช่วยให้สามารถพัฒนาและทดสอบแต่ละส่วนประกอบได้อย่างอิสระ และช่วยให้ AI backend ที่ต้องใช้ทรัพยากรการคำนวณสูงสามารถปรับขนาดหรืออัปเกรดได้โดยไม่ต้องแก้ไขส่วนประสบการณ์ผู้ใช้ใน Frontend

ส่วน Frontend พัฒนาด้วย React ซึ่งเป็น Component-based JavaScript framework ที่สร้างบน HTML \cite{sharma2018html,mozilla2025html} โดยทำหน้าที่เป็นอินเทอร์เฟซสำหรับอัปโหลดรูปภาพ ตัวเลือกการจับภาพจากกล้องแบบเรียลไทม์ และ Dashboard สำหรับแสดงผลลัพธ์แบบโต้ตอบ ส่วน Backend พัฒนาด้วยภาษา Python \cite{siva2023python} โดยใช้ FastAPI \cite{tioman2024fastapi} ซึ่งเป็น Asynchronous web framework ประสิทธิภาพสูงที่ให้บริการ RESTful API เพื่อให้ React frontend เรียกใช้งาน
เมื่อได้รับรูปภาพผ่าน API แล้ว ส่วน Backend จะรัน AI inference pipeline แบบเต็มรูปแบบและส่งคืนผลลัพธ์ในรูปแบบ JSON ที่มีโครงสร้างประกอบด้วยผลการตรวจจับทั้งหมด, การวัดขนาดเส้นผ่านศูนย์กลาง, ป้ายกำกับจาก OCR และการจำแนกความไวต่อยา (Susceptibility classification) นอกจากนี้ยังมีฐานข้อมูลเชิงสัมพันธ์ที่ใช้ SQLite สำหรับจัดเก็บประวัติการวิเคราะห์, บันทึก Session ของผู้ใช้ และตาราง Breakpoint ของ CLSI 2025 และ EUCAST 2023 เพื่อให้สามารถตรวจสอบย้อนหลังและค้นหาเกณฑ์ Breakpoint ได้อย่างรวดเร็วในระหว่างกระบวนการ

Pipeline ของการวิเคราะห์รูปภาพโดยรวมประกอบด้วย 5 ขั้นตอนตามลำดับ ซึ่งแต่ละขั้นตอนจะอธิบายรายละเอียดในหัวข้อถัดไป:
\begin{enumerate}
  \item \textbf{การรับภาพและการประมวลผลเบื้องต้น (Image acquisition and preprocessing):} รูปภาพที่นำเข้าจะได้รับการตรวจสอบความถูกต้อง, ปรับขนาดตามความละเอียดที่โมเดลต้องการ และประมวลผลเบื้องต้นตามข้อกำหนดของโมเดลการตรวจจับที่ใช้งานอยู่
  \item \textbf{การตรวจจับแผ่นยาปฏิชีวนะและการระบุชื่อด้วย OCR (Antibiotic disk detection and OCR labeling):} โมเดลการตรวจจับจะระบุตำแหน่งของแผ่นยาปฏิชีวนะแต่ละแผ่น จากนั้น OCR pipeline จะอ่านรหัสตัวอักษรและตัวเลขจากภาพที่ครอบตัด (Crop) ของแต่ละแผ่นเพื่อระบุชื่อยาปฏิชีวนะ
  \item \textbf{การตรวจจับ ZOI และการวัดเส้นผ่านศูนย์กลาง (ZOI detection and diameter measurement):} โมเดลการตรวจจับจะระบุตำแหน่งของบริเวณยับยั้งเชื้อ (Zone of inhibition) แต่ละบริเวณเป็น Bounding box จากนั้นขนาดของ Bounding box จะถูกแปลงเป็นค่าประมาณของเส้นผ่านศูนย์กลางในหน่วยพิกเซล
  \item \textbf{การสอบเทียบจากพิกเซลเป็นมิลลิเมตร (Pixel-to-millimeter calibration):} การวัดเส้นผ่านศูนย์กลางในหน่วยพิกเซลจะถูกแปลงเป็นหน่วยมิลลิเมตรทางกายภาพ โดยใช้ขนาดเส้นผ่านศูนย์กลาง 6\,มม. ที่ทราบค่าของแผ่นยาปฏิชีวนะเป็นวัตถุอ้างอิงภายในภาพ
  \item \textbf{การค้นหา Breakpoint และการจำแนกความไวต่อยา (Breakpoint lookup and susceptibility classification):} เส้นผ่านศูนย์กลาง ZOI ที่วัดได้ในหน่วยมิลลิเมตรจะถูกนำไปเปรียบเทียบกับ Breakpoint ของ CLSI หรือ EUCAST สำหรับยาปฏิชีวนะที่ระบุได้และชนิดของเชื้อที่ผู้ใช้กำหนด จากนั้นจะจำแนกเชื้อเป็น ไวต่อยา (Susceptible - S), กึ่งไว (Intermediate - I) หรือ ดื้อยา (Resistant - R)
\end{enumerate}

\begin{figure}[H]
  \centering
  \safeimg[width=\textwidth]{figures/pipeline_flowchart.png}{End-to-end analysis pipeline}
  \caption{Pipeline การวิเคราะห์รูปภาพแบบ End-to-end ของระบบ AST อัตโนมัติ ตั้งแต่การนำเข้าภาพผ่านการตรวจจับ, OCR, การสอบเทียบ และการจำแนกความไวต่อยาตาม Breakpoint}
  \label{fig:pipeline}
\end{figure}

\begin{figure}[H]
  \centering
  \safeimg[width=0.92\textwidth]{figures/system_architecture.png}{System Architecture Diagram}
  \caption{สถาปัตยกรรมระบบโดยรวมของเว็บแอปพลิเคชัน AST อัตโนมัติ แสดงให้เห็นส่วน React frontend, FastAPI backend, AI inference pipeline และชั้นฐานข้อมูล SQLite}
  \label{fig:sys_arch}
\end{figure}

\begin{figure}[H]
  \centering
  \safeimg[width=0.82\textwidth]{figures/_review/use_case_diagram.png}{Use case diagram}
  \caption{แผนภาพ Use Case แสดงการโต้ตอบระหว่างผู้ใช้งาน (บุคลากรทางการแพทย์)
    กับระบบวิเคราะห์ AST อัตโนมัติ ครอบคลุมการจัดการบัญชี การวิเคราะห์ภาพ และการดูผลลัพธ์}
  \label{fig:use_case}
\end{figure}

\section{การรวบรวมและการเตรียมชุดข้อมูล (Dataset Collection and Preparation)}

การสร้างชุดข้อมูลการฝึกฝนที่มีคุณภาพสูงและมีความหลากหลายเป็นสิ่งจำเป็นสำหรับประสิทธิภาพการทำงานทั่วไป (Generalization) ของโมเดล Deep Learning สำหรับโดเมน AST การสร้างชุดข้อมูลมีความท้าทายอย่างยิ่งเนื่องจากรูปภาพจานเพาะเชื้อทางคลินิกจริงมักเป็นข้อมูลส่วนบุคคล (เป็นของโรงพยาบาลหรือห้องปฏิบัติการอ้างอิง) และชุดข้อมูลรูปภาพ AST ที่มีการระบุป้ายกำกับ (Annotated) และเผยแพร่สู่สาธารณะมีจำนวนจำกัด กลยุทธ์ชุดข้อมูลที่ใช้ในโปรเจกต์นี้จึงผสมผสานสามแหล่งที่ส่งเสริมกัน ได้แก่ รูปภาพจริงจากห้องปฏิบัติการ, รูปภาพจริงที่ผ่านการเพิ่มข้อมูล (Augmented) และรูปภาพสังเคราะห์ (Synthetic images) ที่สร้างขึ้นโดยใช้โปรแกรม

\subsection{การรวบรวมรูปภาพจริง (Real Image Collection)}

รูปภาพถ่ายจานเพาะเชื้อ AST จริงทั้งหมด 287 ภาพถูกรวบรวมจากสามแหล่ง การแบ่งปันข้อมูลการวิจัยด้านชีวการแพทย์แบบเปิดผ่านคลังข้อมูลเช่น DRYAD เป็นหลักการที่มีมาอย่างยาวนานในการปฏิบัติงานทางวิทยาศาสตร์ \cite{iom1996} แหล่งแรกคือคลังข้อมูล DRYAD \cite{giske2024dryad} ที่เปิดเผยต่อสาธารณะ ซึ่งมีรูปภาพจานเพาะเชื้อแบบ Disk diffusion ที่ระบุป้ายกำกับแล้วจำนวน 225 ภาพ ซึ่งถ่ายในสภาพแวดล้อมห้องปฏิบัติการที่หลากหลาย แหล่งที่สองคือชุดข้อมูล CDS (Clinical and Demographic Study) ซึ่งมีจำนวน 50 ภาพ และแหล่งที่สามคือชุดข้อมูล CRA (Chulabhorn Royal Academy หรือ ราชวิทยาลัยจุฬาภรณ์) ซึ่งมีจำนวน 32 ภาพ เมื่อรวมทั้งสามแหล่งนี้เข้าด้วยกันจึงได้รูปภาพจริงจำนวน 287 ภาพที่ใช้ในโปรเจกต์นี้

รูปภาพเหล่านี้ครอบคลุมเชื้อแบคทีเรียหลายสายพันธุ์, จำนวนแผ่นยาปฏิชีวนะที่แตกต่างกันต่อจาน (ตั้งแต่ 1 ถึง 10 แผ่น) และลักษณะการซ้อนทับกันของบริเวณยับยั้งเชื้อที่หลากหลาย ซึ่งแสดงถึงความหลากหลายของสถานการณ์ทางคลินิกที่ระบบต้องจัดการ ชุดข้อมูลจากแหล่งภายนอกช่วยเพิ่มความแปรผันในเรื่องของแสง, รุ่นของกล้อง, ยี่ห้อของวุ้น (Agar) และขนาดของจานเพาะเชื้อ ซึ่งช่วยปรับปรุงความสามารถในการประมวลผลทั่วไปของโมเดลที่ฝึกฝนด้วยชุดข้อมูลรวม

\subsection{การเพิ่มข้อมูล (Data Augmentation)}

เพื่อขยายชุดข้อมูลการฝึกฝนจาก 287 ภาพ โดยยังคงรักษาความถูกต้องของป้ายกำกับต้นฉบับไว้ จึงได้มีการประยุกต์ใช้การดำเนินการเพิ่มข้อมูลทางเรขาคณิตและเชิงแสงกับรูปภาพจริงโดยใช้แพลตฟอร์ม Roboflow การแปลงข้อมูลรวมถึงการพลิกในแนวตั้งและแนวนอน (Horizontal and vertical flipping), การหมุนแบบสุ่มภายในช่วง $\pm$15$^\circ$, การเปลี่ยนแปลงความสว่าง ($\pm$25\%), การเปลี่ยนแปลงความเปรียบต่าง ($\pm$20\%) และการเพิ่มสัญญาณรบกวน Gaussian (Gaussian noise ที่มีค่าเบี่ยงเบนมาตรฐานสูงสุด 5\% ของค่าความเข้มพิกเซลสูงสุด) รูปภาพที่เพิ่มข้อมูลแต่ละภาพถูกสร้างขึ้นโดยการสุ่มตัวอย่างจากการแปลงเหล่านี้ โดยที่พิกัดของ Bounding box จะถูกคำนวณใหม่โดยอัตโนมัติเพื่อให้ตรงกับเรขาคณิตของรูปภาพที่เปลี่ยนไป กลยุทธ์การเพิ่มข้อมูลนี้ช่วยเพิ่มชุดข้อมูลการฝึกฝนจาก 287 เป็น 543 ภาพ ซึ่งเพิ่มจำนวนตัวอย่างการฝึกฝนที่ไม่ซ้ำกันจากภาพถ่ายจานเพาะเชื้อจริงประมาณสองเท่า

สิ่งสำคัญคือต้องทราบว่าการเพิ่มข้อมูลถูกนำไปใช้เฉพาะกับส่วนของการฝึกฝน (Training partition) ของรูปภาพจริงเท่านั้น ส่วนการตรวจสอบ (Validation) และการทดสอบ (Test partition) จะถูกกันออกจากการเพิ่มข้อมูลเพื่อให้แน่ใจว่าตัวชี้วัดการประเมินผลจะสะท้อนถึงประสิทธิภาพของโมเดลกับข้อมูลนำเข้าที่ไม่ได้ดัดแปลงและมีความสมจริง นอกจากนี้ การเพิ่มข้อมูลยังถูกจำกัดไว้เฉพาะการแปลงที่สอดคล้องกับความแปรปรวนที่อาจเกิดขึ้นจริงในการถ่ายภาพในห้องปฏิบัติการ โดยหลีกเลี่ยงการแปลงที่รุนแรงเกินไป (เช่น การบิดเบือนมุมมองอย่างหนัก หรือการสลับสี) ซึ่งจะสร้างภาพที่ไม่สมจริงและส่งผลเสียต่อการฝึกฝนโมเดล

\subsection{การสร้างข้อมูลสังเคราะห์ (Synthetic Data Generation)}

รูปภาพสังเคราะห์อีก 2,000 ภาพถูกสร้างขึ้นโดยใช้โปรแกรมเพื่อเพิ่มความหลากหลายของชุดข้อมูลและเพื่อแก้ปัญหาการขาดแคลนรูปภาพจริงที่มีลักษณะการซ้อนทับของบริเวณยับยั้งเฉพาะจุดและรูปร่างบริเวณยับยั้งที่ผิดปกติ รูปภาพจานเพาะเชื้อสังเคราะห์ถูกเรนเดอร์โดยใช้ Pipeline เชิงกระบวนการ (Procedural pipeline) ที่พัฒนาด้วย Python โดยใช้ไลบรารี Pillow และ OpenCV ในการประกอบภาพ รูปภาพสังเคราะห์แต่ละภาพถูกออกแบบมาให้มีลักษณะทางสายตาใกล้เคียงกับจาน Mueller-Hinton agar จริง โดยมีพื้นหลังวุ้นที่มีลวดลาย (Texture) ซึ่งสร้างจากการเรียงต่อกันและผสมผสานลวดลายตัวอย่าง, วัตถุแผ่นยาปฏิชีวนะทรงกลมที่วางในตำแหน่งที่กำหนด และบริเวณยับยั้งเชื้อทรงกลมที่เรนเดอร์ด้วยขอบฟุ้งแบบ Gaussian เพื่อเลียนแบบขอบเขตการเจริญเติบโตที่ค่อยๆ จางหายไปตามที่เห็นในภาพถ่ายจริง

\begin{figure}[H]
  \centering
  \safeimg[width=0.58\textwidth]{figures/_review/synthetic_plate_example.png}{Synthetic plate example}
  \caption{ตัวอย่างภาพ Synthetic ที่สร้างขึ้นด้วยโปรแกรม แสดงจานเพาะเชื้อจำลองพร้อมแผ่นยาปฏิชีวนะ
    และวงยับยั้งที่มีขนาดหลากหลาย ใช้สำหรับเพิ่มปริมาณและความหลากหลายของชุดข้อมูลฝึกฝน}
  \label{fig:synthetic_example}
\end{figure}

รูปภาพสังเคราะห์ 2,000 ภาพถูกแบ่งออกเป็นสองประเภทในจำนวนที่เท่ากัน ประเภทแรกจำนวน 1,000 ภาพใช้ตำแหน่งแผ่นยาปฏิชีวนะแบบคงที่ตามรูปแบบจานเพาะเชื้อที่แนะนำโดย CLSI สำหรับการวางแผ่นยา 6 แผ่นและ 8 แผ่น เพื่อจำลองรูปแบบมาตรฐานที่เป็นระบบในขั้นตอนการทำงานของห้องปฏิบัติการทั่วไป ประเภทที่สองจำนวน 1,000 ภาพใช้การวางแผ่นยาแบบสุ่มภายในขอบเขตของจาน เพื่อจำลองการจัดวางที่เป็นระบบน้อยกว่าซึ่งพบได้บ่อยในห้องปฏิบัติการวิจัยและการเรียนการสอน ทั้งสองประเภทมีการสุ่มขนาดเส้นผ่านศูนย์กลางของบริเวณยับยั้งจากช่วงค่าทางคลินิกที่สมจริง (ตั้งแต่ประมาณ 10\,มม. ถึง 40\,มม.) โดยจงใจเพิ่มสัดส่วนของสถานการณ์ที่บริเวณยับยั้งมีการซ้อนทับกัน (ซึ่งบริเวณยับยั้งของแผ่นยาที่อยู่ใกล้กันบางส่วนตัดกัน) เพื่อให้แน่ใจว่าโมเดลได้เรียนรู้กรณีที่ท้าทายแต่มีความสำคัญทางคลินิกนี้ในระหว่างการฝึกฝน ลวดลายพื้นหลัง, การจำลองการไล่ระดับแสง และรูปแบบการเจริญเติบโตของโคโลนีสังเคราะห์ก็ถูกทำให้หลากหลายเพื่อเพิ่มความหลากหลายทางสายตา

\subsection{การแบ่งชุดข้อมูลขั้นสุดท้าย (Final Dataset Split)}

\begin{table}[H]
  \centering
  \caption{องค์ประกอบของชุดข้อมูลและการแบ่งส่วนสำหรับการฝึกฝน/การตรวจสอบ/การทดสอบ ที่ใช้ในการฝึกฝนและประเมินผลโมเดล}
  \label{tab:dataset_split}
  \begin{tabular}{lcccc}
    \toprule
    \textbf{แหล่งข้อมูล (Source)} & \textbf{ทั้งหมด} & \textbf{ฝึกฝน (Train)} & \textbf{ตรวจสอบ (Validation)} & \textbf{ทดสอบ (Test)} \\
    \midrule
    DRYAD                              & 225   & ---   & ---   & ---   \\
    CDS                                & 50    & ---   & ---   & ---   \\
    CRA (ราชวิทยาลัยจุฬาภรณ์)           & 32    & ---   & ---   & ---   \\
    \textit{รวมรูปภาพจริง (ต้นฉบับ)}     & 287   & ---   & ---   & ---   \\
    \textit{รวมรูปภาพจริง (เพิ่มข้อมูล)}  & 543   & 433   & 55    & 55    \\
    ข้อมูลสังเคราะห์ (ตำแหน่งคงที่)       & 1,000 & \multicolumn{3}{c}{ใช้ในการฝึกฝนเท่านั้น} \\
    ข้อมูลสังเคราะห์ (ตำแหน่งไม่คงที่)     & 1,000 & \multicolumn{3}{c}{ใช้ในการฝึกฝนเท่านั้น} \\
    \bottomrule
  \end{tabular}
\end{table}

รูปภาพสังเคราะห์ถูกนำมาใช้เพื่อการฝึกฝนเท่านั้น เนื่องจากถูกสร้างขึ้นมาเพื่อเสริมการเรียนรู้ของโมเดลในรูปแบบที่หลากหลาย มากกว่าที่จะใช้ประเมินประสิทธิภาพการทำงานทั่วไป ตัวชี้วัดประสิทธิภาพทั้งหมดที่รายงานในบทที่ 4 คำนวณจากชุดทดสอบรูปภาพจริง (55 ภาพ) เพื่อให้แน่ใจว่าตัวเลขประสิทธิภาพที่รายงานสะท้อนถึงพฤติกรรมของโมเดลกับภาพถ่ายในห้องปฏิบัติการที่เป็นของจริงและไม่ได้ถูกตกแต่ง

\section{การระบุป้ายกำกับ (Annotation)}

ป้ายกำกับ Ground-truth ถูกสร้างขึ้นในรูปแบบ Pascal VOC XML โดยระบุพิกัด Bounding box สำหรับแต่ละบริเวณยับยั้งเชื้อและแต่ละแผ่นยาปฏิชีวนะที่ปรากฏในรูปภาพ รูปแบบการระบุป้ายกำตอนนี้จะเข้ารหัสค่าสี่ค่าต่อวัตถุ ได้แก่ พิกัดพิกเซลของมุมบนซ้ายและมุมล่างขวาของ Bounding box พร้อมกับป้ายกำกับคลาส (\texttt{antibiotic} สำหรับแผ่นยา และ \texttt{disk-zone} สำหรับบริเวณยับยั้งเชื้อ)

การระบุป้ายกำกับเบื้องต้นทำโดยทีมงานโปรเจกต์โดยใช้ LabelImg ซึ่งเป็นเครื่องมือระบุป้ายกำกับรูปภาพแบบกราฟิกที่เป็นโอเพนซอร์ส เพื่อให้แน่ใจว่ามีความถูกต้องทางคลินิก การวาง Bounding box ทั้งหมดสำหรับคลาสบริเวณยับยั้งเชื้อได้รับการตรวจสอบและแก้ไขโดยนักจุลชีววิทยาที่ผ่านการฝึกฝนมาแล้ว ซึ่งเป็นผู้ตรวจสอบป้ายกำกับแต่ละรายการเปรียบเทียบกับการวัดจานจริงโดยใช้ Calipers ขั้นตอนการตรวจสอบสองชั้นนี้เป็นสิ่งจำเป็นสำหรับการสร้าง Ground-truth ที่เชื่อถือได้ เนื่องจากความหมายของขอบเขตบริเวณยับยั้งต้องใช้ความเชี่ยวชาญเฉพาะทาง: ขอบที่ถูกต้องของบริเวณยับยั้งถูกกำหนดให้เป็นจุดที่การยับยั้งการเจริญเติบโตแบบ Confluent growth สิ้นสุดลงโดยสมบูรณ์ (ไม่ใช่บริเวณที่มีการเกิด Micro-colony ซึ่ง CLSI แนะนำให้เจ้าหน้าที่ละเว้น) จากนั้นป้ายกำกับจะถูกนำเข้าสู่ Roboflow เพื่อจัดการชุดข้อมูล, กำหนด Pipeline การเพิ่มข้อมูล และส่งออกในรูปแบบที่เฟรมเวิร์กการฝึกฝนของแต่ละโมเดลต้องการ (Pascal VOC XML สำหรับ Faster R-CNN และ YOLO TXT สำหรับ YOLOv11n)

\begin{figure}[H]
  \centering
  \safeimg[width=0.82\textwidth]{figures/_review/annotation_roboflow.png}{Roboflow annotation screenshot}
  \caption{หน้าจอการระบุป้ายกำกับ (Annotation) บนแพลตฟอร์ม Roboflow แสดง Bounding box
    สำหรับคลาส \texttt{antibiotic} (ดิสก์ยาปฏิชีวนะ สีแดง) และ \texttt{disk-zone}
    (วงยับยั้งเชื้อ สีม่วง) พร้อมสถิติจำนวนตัวอย่างในแต่ละคลาส}
  \label{fig:annotation}
\end{figure}

\section{การฝึกฝนโมเดล (Model Training)}

\subsection{Faster R-CNN}

Faster R-CNN ถูกพัฒนาโดยใช้ไลบรารี \texttt{torchvision} ของ PyTorch พร้อมกับ Backbone แบบ ResNet-50 ที่ผ่านการฝึกฝนเบื้องต้น (Pretrained) บนชุดข้อมูล ImageNet \cite{he2016deep} การใช้ Backbone ที่ผ่านการฝึกฝนมาแล้วเป็นสิ่งสำคัญมากสำหรับประสิทธิภาพของชุดข้อมูลภาพทางการแพทย์ที่มีขนาดเท่านี้: การฝึกฝนเบื้องต้นบน ImageNet ช่วยให้ Backbone มีตัวแทนคุณลักษณะ (Feature representations) ที่แยกแยะโครงสร้างรูปภาพระดับต่ำ (ขอบ, พื้นผิว, การไล่ระดับสี) ซึ่งสามารถถ่ายโอนมายังโดเมน AST ได้อย่างมีประสิทธิภาพ ช่วยเร่งการลู่เข้าหาคำตอบ (Convergence) และปรับปรุงความแม่นยำขั้นสุดท้ายได้ดีกว่าการเริ่มฝึกฝนจากการสุ่มค่าเริ่มต้น

ส่วนหัวของการตรวจจับ (Detection head ทั้ง RPN และ Classification/regression head) ถูกกำหนดค่าเริ่มต้นใหม่และฝึกฝนตั้งแต่ต้น การฝึกฝนใช้อัลกอริทึมการเพิ่มประสิทธิภาพ Stochastic Gradient Descent (SGD) โดยมี Learning rate เท่ากับ $0.005$, Momentum เท่ากับ $0.9$ และ Weight decay เท่ากับ $5 \times 10^{-4}$ (คือ \texttt{0.0005}) เพื่อรักษาสมดุลของโมเดลและป้องกันการ Overfitting กับชุดข้อมูลรูปภาพจริงที่มีอยู่อย่างจำกัด ค่า Learning rate จะถูกปรับลดลงโดยใช้ \texttt{StepLR} scheduler โดยมี \texttt{step\_size=3} และ \texttt{gamma=0.1} ซึ่งจะลด Learning rate ลง 10 เท่าทุกๆ 3 Epochs รูปภาพทั้งหมดถูกปรับขนาดเป็น $640 \times 640$ พิกเซลก่อนนำเข้าโมเดล และใช้ Batch size เท่ากับ 8

ก่อนที่จะเริ่มการฝึกฝนแบบเต็มรูปแบบ ได้มีการรันการทดลองฝึกฝนเบื้องต้นบนชุดย่อยของชุดข้อมูลที่ขยายขนาดขึ้นเรื่อยๆ เพื่อศึกษาพฤติกรรมการเรียนรู้ของโมเดลและวินิจฉัยสภาวะ Underfitting และ Overfitting การทดลองเหล่านี้ช่วยในการเลือก Hyperparameter ขั้นสุดท้าย โดยเฉพาะจำนวน Epochs ในการฝึกฝน (10 Epochs) และตารางเวลาของ Learning rate โดยแสดงให้เห็นว่า Epoch ใดที่ Loss ของการตรวจสอบ (Validation loss) มีค่าต่ำสุดก่อนที่จะเริ่มสูงขึ้นอีกครั้ง การฝึกฝนขั้นสุดท้ายใช้เวลาประมาณ 174 นาทีบนฮาร์ดแวร์ที่มีอยู่ ซึ่งสะท้อนถึงต้นทุนการคำนวณของสถาปัตยกรรม Inference แบบสองขั้นตอน

\subsection{YOLOv11n}

YOLOv11n ถูกฝึกฝนโดยใช้เฟรมเวิร์กการฝึกฝนของ Ultralytics ซึ่งให้บริการ Python API ระดับสูงสำหรับการกำหนดค่าเริ่มต้นโมเดล, การดำเนินการลูปการฝึกฝน และการบันทึก Checkpoint \cite{wang2024yolov11} ตัวเลือกย่อยรุ่น Nano (n) ถูกเลือกเพื่อเน้นความเร็วในการประมวลผล (Inference speed) และความสามารถในการนำไปใช้งานบนระบบที่ไม่มี GPU แยก ในขณะที่ยังคงรักษาสมรรถนะของโมเดลเพียงพอที่จะเรียนรู้คุณลักษณะทางสายตาที่ค่อนข้างจำกัดของรูปภาพจานเพาะเชื้อ AST (คลาสวัตถุสองคลาสที่มีรูปร่างที่ชัดเจน: แผ่นยาทรงกลม และบริเวณยับยั้งทรงกลม)

น้ำหนักของโมเดลถูกกำหนดค่าเริ่มต้นจาก Checkpoint ที่ผ่านการฝึกฝนบน COCO ของ Ultralytics สำหรับ YOLOv11n การเรียนรู้แบบถ่ายโอน (Transfer learning) จาก COCO ซึ่งเป็นฐานข้อมูลการตรวจจับวัตถุในธรรมชาติขนาดใหญ่ที่มีวัตถุ 80 ประเภท ช่วยให้โมเดลมีความสามารถในการตรวจจับวัตถุทั่วไป (การดึงคุณลักษณะ, การระงับค่าที่ไม่ใช่ค่าสูงสุดหรือ Non-maximum suppression, การพยากรณ์แบบ Anchor-free) ซึ่งจะถูกปรับจูน (Fine-tuning) ให้เข้ากับโดเมนเฉพาะอย่าง AST การฝึกฝนแบบ Fine-tuning ถูกรันเป็นจำนวน \textbf{200 Epochs} บนการแบ่งชุดข้อมูลเดียวกันกับ Faster R-CNN โดยรูปภาพถูกปรับขนาดเป็น $640 \times 640$ พิกเซล และใช้ Batch size เท่ากับ 8 โดยได้ปิดฟังก์ชันการเพิ่มข้อมูลในตัวของ Ultralytics (\texttt{augment=False}) เพื่อให้โมเดลได้รับการฝึกฝนด้วยรูปภาพที่สะอาดและไม่มีการเพิ่มข้อมูล ซึ่งสอดคล้องกับกลยุทธ์การเพิ่มข้อมูลที่ใช้ในขั้นตอนการเตรียมชุดข้อมูล การฝึกฝนด้วยชุดข้อมูลเต็มรูปแบบในการกำหนดค่านี้ใช้เวลาประมาณ 200 นาที

YOLOv11n ตรวจจับวัตถุสองคลาส ได้แก่ \texttt{antibiotic} (แผ่นยาที่มีตัวยา ซึ่งทำหน้าที่เป็นเป้าหมายของ OCR และเป็นตัวอ้างอิงในการสอบเทียบในภาพ) และ \texttt{disk-zone} (บริเวณยับยั้งโดยรอบ ซึ่งใช้สำหรับวัดเส้นผ่านศูนย์กลาง ZOI) ในระหว่างการประมวลผล โมเดลจะส่งคืน Bounding boxes สำหรับวัตถุที่ตรวจจับได้ทั้งหมดจากทั้งสองคลาส ซึ่งจะถูกจับคู่ตามความสัมพันธ์เชิงพื้นที่ (การจับคู่แผ่นยาที่ใกล้ที่สุดกับบริเวณยับยั้ง) เพื่อเชื่อมโยงแต่ละบริเวณยับยั้งเข้ากับแผ่นยาปฏิชีวนะที่เกี่ยวข้อง

\subsection{Hough Circle Transform}

Hough Circle Transform ถูกนำมาใช้เป็นเกณฑ์มาตรฐานแบบคลาสสิกโดยใช้ฟังก์ชัน \texttt{HoughCircles} ของ OpenCV ซึ่งประยุกต์ใช้ Generalized Hough transform สำหรับการตรวจจับวงกลมโดยใช้การลงคะแนนตามความชัน (Gradient-based voting) ก่อนที่จะประยุกต์ใช้การแปลงนี้ รูปภาพแต่ละภาพจะถูกแปลงเป็นสีเทาและทำให้เรียบขึ้นด้วย Gaussian filter โดยใช้ขนาด Kernel $17 \times 17$ พิกเซล และค่า $\sigma = 0$ (เพื่อให้ OpenCV คำนวณค่าเบี่ยงเบนมาตรฐานโดยอัตโนมัติจากขนาด Kernel) เพื่อลดสัญญาณรบกวนที่อาจก่อให้เกิดการตอบสนองต่อความชันที่ผิดพลาด

ฟังก์ชัน \texttt{HoughCircles} รับพารามิเตอร์หลักสี่ตัว ได้แก่ \texttt{dp} (อัตราส่วนความละเอียดผกผันของ Accumulator ต่อภาพนำเข้า ซึ่งควบคุมความละเอียดของ Accumulator ใน Parameter-space), \texttt{minDist} (ระยะห่างขั้นต่ำที่อนุญาตระหว่างจุดศูนย์กลางของวงกลมที่ตรวจจับได้ เพื่อป้องกันการตรวจจับซ้ำซ้อนในบริเวณเดียวกัน), \texttt{param1} (เกณฑ์การตรวจจับขอบ Canny ตัวบน โดยเกณฑ์ตัวล่างจะถูกตั้งเป็นครึ่งหนึ่งของค่านี้โดยอัตโนมัติ) และ \texttt{param2} (เกณฑ์ Accumulator สำหรับการตรวจจับจุดศูนย์กลาง ซึ่งควบคุมจำนวนการลงคะแนนขั้นต่ำที่จำเป็นในการประกาศวงกลมที่ตรวจจับได้) ค่าที่เหมาะสมที่สุดสำหรับพารามิเตอร์ทั้งสี่นี้ได้มาจากการค้นหาแบบตาราง (Grid search) อย่างละเอียดทั่วทั้งชุดข้อมูล เพื่อลดค่าความผิดพลาดสัมบูรณ์เฉลี่ย (Mean absolute error) ระหว่างรัศมีของวงกลมที่ตรวจจับได้กับค่า Ground-truth ค่าสุดท้ายที่เลือกคือ \texttt{dp=1.5}, \texttt{minDist=150}, \texttt{param1=100} และ \texttt{param2=50} สิ่งสำคัญคือตัวเลขประสิทธิภาพสุดท้ายที่รายงานนี้เป็นกรณีที่ดีที่สุดสำหรับ Hough Transform เนื่องจากค่าพารามิเตอร์ถูกปรับจูนโดยใช้ชุดข้อมูลทั้งหมด ซึ่งน่าจะรวมถึงชุดทดสอบด้วย ดังนั้นประสิทธิภาพในการทำงานทั่วไปอาจต่ำกว่าตัวเลขที่รายงาน

\section{การวัดเส้นผ่านศูนย์กลาง ZOI และการสอบเทียบ (ZOI Diameter Measurement and Calibration)}

หลังจากตรวจจับบริเวณยับยั้งในฐานะ Bounding boxes (สำหรับโมเดล Deep Learning ทั้งสอง) หรือเป็นวงกลมที่มีพารามิเตอร์รัศมีที่ชัดเจน (สำหรับ Hough Transform) ผลลัพธ์ดิบจะต้องถูกแปลงเป็นหน่วยวัดเส้นผ่านศูนย์กลาง ZOI ทางกายภาพในหน่วยมิลลิเมตรก่อนที่จะนำไปเปรียบเทียบกับ Breakpoints ของ CLSI หรือ EUCAST

สำหรับการตรวจจับที่ใช้ Bounding box (Faster R-CNN และ YOLOv11n) เส้นผ่านศูนย์กลางในหน่วยพิกเซลของบริเวณที่ตรวจจับได้จะถูกประมาณจากขนาดของ Bounding box โดยใช้การประมาณสองแบบที่นำมาเปรียบเทียบและหาค่าเฉลี่ย:
\begin{itemize}
  \item \textbf{อิงตามเส้นผ่านศูนย์กลาง (Diameter-based):} ค่าเฉลี่ยของความกว้าง $w$ และความสูง $h$ ของ Bounding box ในหน่วยพิกเซล คือ $\hat{d}_{\mathrm{diam}} = (w + h)/2$ โดยสมมติว่าบริเวณนั้นเป็นวงกลมโดยประมาณ
  \item \textbf{อิงตามพื้นที่ (Area-based):} เส้นผ่านศูนย์กลางของวงกลมที่มีพื้นที่เท่ากันซึ่งคำนวณจากพื้นที่ Bounding box คือ $\hat{d}_{\mathrm{area}} = 2\sqrt{(w \cdot h)/\pi}$ ซึ่งมีความทนทานมากกว่าสำหรับบริเวณที่มีลักษณะเป็นวงรีเล็กน้อยที่อัตราส่วนภาพของ Bounding box เบี่ยงเบนไปจากหนึ่ง
\end{itemize}

การแปลงจากพิกเซลเป็นมิลลิเมตรดำเนินการโดยใช้แผ่นยาปฏิชีวนะเป็นตัวอ้างอิงในการสอบเทียบภายในภาพ เนื่องจากเส้นผ่านศูนย์กลางทางกายภาพของแผ่นยาปฏิชีวนะมาตรฐานตามข้อกำหนดของ CLSI คือ 6\,มม. พอดี และแผ่นยาก็ถูกตรวจจับเป็น Bounding box โดยโมเดลเช่นกัน ดังนั้นปัจจัยการสอบเทียบ (Calibration factor ในหน่วยมิลลิเมตรต่อพิกเซล) จึงสามารถคำนวณได้สำหรับแต่ละภาพแยกกันดังนี้ $k = 6\,\mathrm{mm} / \hat{d}_{\mathrm{disk}}$ โดยที่ $\hat{d}_{\mathrm{disk}}$ คือค่าประมาณเส้นผ่านศูนย์กลางในหน่วยพิกเซลของ Bounding box ของแผ่นยาปฏิชีวนะที่เกี่ยวข้อง การใช้ปัจจัยการสอบเทียบต่อภาพนี้ช่วยขจัดผลกระทบของระยะห่างของกล้องและระดับการซูมที่แตกต่างกัน ทำให้มั่นใจได้ว่าการวัดมีความแม่นยำแม้จะมีความแปรปรวนของการถ่ายภาพ โดยไม่จำเป็นต้องใช้ไม้บรรทัดจริงในมุมมองของกล้อง

การปรับปรุงการสอบเทียบเป็นกระบวนการที่ทำซ้ำตลอดทั้งโปรเจกต์ การสอบเทียบเบื้องต้นโดยใช้การประมาณเส้นผ่านศูนย์กลางอิงตามพื้นที่เพียงอย่างเดียวส่งผลให้เกิดความคลาดเคลื่อนเชิงระบบในทิศทางบวกประมาณ $+4$\,มม. (ระบบประเมินเส้นผ่านศูนย์กลางบริเวณยับยั้งสูงเกินจริงอย่างต่อเนื่อง) การวิเคราะห์สาเหตุของความคลาดเคลื่อนนี้พบว่า Bounding box ที่พยากรณ์โดยโมเดลมักจะกว้างกว่าขอบเขตบริเวณยับยั้งจริงเล็กน้อยเนื่องจากขอบของบริเวณยับยั้งมีความฟุ้งและไล่ระดับสี การผสมผสานการประมาณแบบอิงตามพื้นที่และอิงตามเส้นผ่านศูนย์กลางเข้าด้วยกัน และการประยุกต์ใช้ปัจจัยการแก้ไขเชิงเส้นที่เรียนรู้มา (Learned linear correction factor) ช่วยลดความคลาดเคลื่อนเชิงระบบจากประมาณ $+4$\,มม. เหลือประมาณ $+2.1$\,มม. ซึ่งถือเป็นการปรับปรุงความแม่นยำในการวัดอย่างมีนัยสำคัญ

\section{OCR Pipeline สำหรับการระบุชื่อยาปฏิชีวนะ (OCR Pipeline for Antibiotic Label Recognition)}

OCR pipeline มีหน้าที่อ่านรหัสตัวอักษรและตัวเลขของยาปฏิชีวนะจากแผ่นยาแต่ละแผ่นที่ตรวจจับได้ และจับคู่กับชื่อยาปฏิชีวนะที่ได้รับการยอมรับในฐานข้อมูล Breakpoint ระบบนี้ได้รับการออกแบบด้วยกลยุทธ์แบบคู่ (Dual-strategy approach) เพื่อรับมือกับความหลากหลายของคุณภาพรูปภาพแผ่นยาที่พบในสภาพแวดล้อมห้องปฏิบัติการและยี่ห้อของแผ่นยาที่แตกต่างกัน

\subsection{กลยุทธ์การครอบตัดแผ่นยาและการประมวลผลเบื้องต้น (Disk Crop and Preprocessing Strategy)}

แต่ละบริเวณแผ่นยาปฏิชีวนะจะถูกครอบตัด (Crop) ออกจากรูปภาพจานเพาะเชื้อตัวเต็มก่อน โดยใช้ Bounding box ที่พยากรณ์โดยโมเดล YOLO สำหรับคลาส \texttt{antibiotic} โดยมีการเพิ่มขอบ (Padding) เล็กน้อยรอบเขตการครอบตัดเพื่อให้แน่ใจว่าได้พื้นที่ผิวของแผ่นยาครบถ้วน แม้ว่า Bounding box จะประเมินขอบเขตของแผ่นยาต่ำไปเล็กน้อยก็ตาม

กลยุทธ์การประมวลผลเบื้องต้นแบบคู่จะประยุกต์ใช้ OCR ในสอง Pipeline แยกกัน และเลือกผลลัพธ์ที่มีคะแนนความมั่นใจ (Confidence score) สูงกว่า:
\begin{enumerate}
  \item \textbf{เส้นทางรูปภาพดิบ (Raw image path):} ประยุกต์ใช้ EasyOCR \cite{jaidedai2024easyocr} โดยตรงกับภาพแผ่นยาที่ครอบตัดแบบสีเทาตามค่าความเข้มพิกเซลเดิม โดยไม่มีการแปลงเพื่อประมวลผลเบื้องต้น เส้นทางนี้ทำงานได้ดีเมื่อภาพแผ่นยามีความเปรียบต่างสูงอยู่แล้วและข้อความตัวหมึกอ่านได้ชัดเจนโดยไม่ต้องปรับปรุง
  \item \textbf{เส้นทาง CLAHE + Otsu (CLAHE + Otsu path):} ภาพแผ่นยาที่ครอบตัดจะถูกประมวลผลด้วย CLAHE (ขนาด Tile $8 \times 8$ พิกเซล, ขีดจำกัดความเปรียบต่าง 2.0) เพื่อปรับการแปรผันของแสงในระดับท้องถิ่นให้เป็นปกติ จากนั้นจะถูกทำเป็นภาพขาวดำ (Binarized) โดยใช้เกณฑ์สากลของ Otsu เพื่อสร้างภาพไบนารีที่มีความเปรียบต่างสูง การดำเนินการทางสัณฐานวิทยา (Morphological dilation or erosion) จะถูกประยุกต์ใช้แบบปรับตัวตามอัตราส่วนพิกเซลข้อความ (เศษส่วนของพิกเซลสีขาวในภาพครอบตัดที่เป็นไบนารี): หากอัตราส่วนต่ำกว่า 0.05 ซึ่งบ่งชี้ว่าตัวอักษรบางหรือจางมาก จะใช้ $3 \times 3$ Dilation kernel เพื่อฟื้นฟูความต่อเนื่องของเส้นตัวอักษร หากอัตราส่วนเกิน 0.45 ซึ่งบ่งชี้ว่าตัวอักษรหนาเกินไปหรือรวมกัน จะใช้ $3 \times 3$ Erosion kernel จากนั้นจึงประยุกต์ใช้ EasyOCR กับภาพไบนารีที่ผ่านการประมวลผลเบื้องต้นนี้
\end{enumerate}

หลังจากได้ข้อความที่เป็นไปได้จากทั้งสอง Pipeline แล้ว OCR pipeline จะดำเนินการปรับทิศทางให้เป็นปกติ (Orientation normalization): โดยการรัน EasyOCR ที่ \textbf{24 ทิศทาง} ของภาพครอบตัด (ตั้งแต่ -180$^\circ$ ถึง +180$^\circ$ โดยก้าวทีละ 15$^\circ$) สำหรับแต่ละ Pipeline ซึ่งจะสร้างสตริงข้อความที่เป็นไปได้หลายรายการ กลยุทธ์การหมุนที่มีความละเอียดสูงนี้ช่วยให้มั่นใจได้ว่าสามารถอ่านป้ายกำกับได้ไม่ว่าแผ่นยาจะวางตัวอย่างไรบนพื้นผิววุ้น ข้อความที่เป็นไปได้ที่มีคะแนนความมั่นใจของ EasyOCR ที่ได้รับการปรับแต่ง (\emph{adjusted} EasyOCR confidence score โดยคำนวณร่วมกับการตรวจสอบรูปแบบด้วย Regex) สูงที่สุดจากทุกทิศทางและทุก Pipeline จะถูกเลือกเป็นผลลัพธ์ OCR ขั้นสุดท้าย

\section{การติดตั้งระบบและโครงสร้างพื้นฐาน (System Deployment and Infrastructure)}

ระบบขั้นสุดท้ายได้รับการเปลี่ยนผ่านจากสภาพแวดล้อมการพัฒนาในเครื่อง (Local) ไปยังการติดตั้งใช้งานในระดับการผลิต (Production-equivalent deployment) ที่เซิร์ฟเวอร์ของ \textbf{ราชวิทยาลัยจุฬาภรณ์ (Chulabhorn Royal Academy หรือ CRA)} โครงสร้างพื้นฐานสำหรับการผลิตนี้ถูกออกแบบมาเพื่อความพร้อมใช้งานสูง (High availability) และการเข้าถึงทางคลินิกที่ปลอดภัย

\subsection{สภาพแวดล้อมเซิร์ฟเวอร์และบริการ Backend (Server Environment and Backend Services)}

บริการ Backend ถูกโฮสต์บนสภาพแวดล้อมเซิร์ฟเวอร์ระบบ Linux ที่ CRA แอปพลิเคชัน FastAPI ให้บริการโดยใช้ \textbf{Uvicorn} ซึ่งเป็น ASGI server ที่รวดเร็วมาก โดยรันเป็นกระบวนการพื้นหลัง (Background process) ที่ได้รับการจัดการ เพื่อให้มั่นใจในเสถียรภาพของระบบ จึงมีสคริปต์จัดการ \texttt{start.sh} โดยเฉพาะที่ทำหน้าที่ประสานงานบริการ, ตรวจสอบกระบวนการ และกู้คืนระบบโดยอัตโนมัติ Backend ถูกกำหนดค่าให้ฟังคำขอ (Listen) บนพอร์ตหมายเลขสูง (6014) เพื่อหลีกเลี่ยงความขัดแย้งกับบริการอื่นๆ ของโรงพยาบาล

\subsection{Reverse Proxy และความปลอดภัย (Reverse Proxy and Security)}

มีการติดตั้ง \textbf{Nginx} เป็น Reverse proxy และ Web server เพื่อจัดการการรับส่งข้อมูลที่เข้ามายังแอปพลิเคชัน Nginx ทำหน้าที่จัดการ SSL/TLS termination เพื่อให้สามารถเข้าถึงระบบผ่าน HTTPS ได้อย่างปลอดภัย โดยมีการกำหนดค่าบล็อก \texttt{location} เฉพาะสำหรับการกำหนดเส้นทาง:
\begin{itemize}
    \item \texttt{/zoneanalyzer/}: ส่งต่อข้อมูลไปยัง React frontend (npx serve บนพอร์ต 6015)
    \item \texttt{/zoneanalyzer/api/}: ส่งต่อคำขอไปยัง FastAPI backend (พอร์ต 6014)
    \item \texttt{/zoneanalyzer/uploaded_images/}: ให้บริการรูปภาพที่แสดงผลลัพธ์โดยตรงจากพื้นที่จัดเก็บของเซิร์ฟเวอร์
\end{itemize}
การกำหนดค่านี้ช่วยให้ใช้จุดเชื่อมต่อโดเมนเดียว (\texttt{https://paniti-jupyter.cra.ac.th/zoneanalyzer/}) สำหรับทั้งแอปพลิเคชัน ช่วยลดความยุ่งยากในการติดตั้งทางคลินิกและการจัดการไฟร์วอลล์

\subsection{การให้บริการ Frontend (Frontend Serving)}

React frontend ได้รับการปรับแต่งเพื่อการผลิตโดยใช้ Vite build pipeline ซึ่งจะสร้างชุดไฟล์ Static assets ที่มีขนาดเล็กและผ่านการทำ Tree-shaken ไฟล์เหล่านี้ถูกให้บริการโดยใช้ \textbf{npx serve} ซึ่งเป็นกลไกการส่งมอบอินเทอร์เฟซผู้ใช้ที่มีน้ำหนักเบาและมีประสิทธิภาพ สภาพแวดล้อม Frontend ถูกกำหนดค่าด้วยไฟล์ \texttt{.env} เฉพาะที่ชี้ไปยัง Endpoint ของ Production API เพื่อให้แน่ใจว่ามีการสื่อสารที่ราบรื่นผ่าน Reverse proxy

\subsection{การประมวลผลหลังการทำงานและการจับคู่ชื่อยาปฏิชีวนะ (Post-Processing and Antibiotic Name Matching)}

ผลลัพธ์ OCR ดิบจะถูกกรองผ่าน Regular expression เพื่อปฏิเสธสตริงที่ไม่น่าจะเป็นไปได้ รูปแบบ Regular expression \texttt{\^{}([A-Za-z]+(\textbackslash{}s[A-Za-z]+)\{0,3\})\textbackslash{}s+(\textbackslash{}d+)\$} จะบังคับใช้รูปแบบที่คาดหวังของชื่อบนแผ่นยาปฏิชีวนะ ได้แก่ โทเค็นตัวอักษรหนึ่งถึงสี่ตัว (ตัวย่อยา) ตามด้วยโทเค็นตัวเลข (ความเข้มข้นในหน่วยไมโครกรัม) สตริงที่ไม่ตรงกับรูปแบบนี้ ซึ่งรวมถึงการหลอนของ OCR (Hallucinations), การอ่านค่าได้เพียงบางส่วน และสตริงที่มีแต่ตัวเลข จะถูกตัดทิ้ง และแผ่นยานั้นจะถูกทำเครื่องหมายว่าเป็น ``ไม่สามารถระบุได้ (Unidentified)'' ในผลลัพธ์

ตัวย่อยาปฏิชีวนะที่จับคู่ได้จะถูกนำไปค้นหาในตารางจับคู่ชื่อยาปฏิชีวนะที่จัดเตรียมไว้ ซึ่งจะจับคู่ตัวย่อมาตรฐานบนแผ่นยา (เช่น ``OX'', ``AMX'', ``CIP'') กับชื่อเต็มของยาปฏิชีวนะและข้อมูลประเภทของยา การจับคู่นี้ช่วยให้การค้นหา Breakpoint สามารถดึงเกณฑ์ของ CLSI หรือ EUCAST ที่ถูกต้องสำหรับยาที่ระบุและชนิดของเชื้อที่ผู้ใช้กำหนดได้

\section{การจำแนกความไวต่อยา (Susceptibility Classification)}

ระบบได้นำตรรกะ Breakpoint ของ CLSI 2025 \cite{clsi2025} และ EUCAST 2023 \cite{eucast2023} มาใช้สำหรับการจำแนกความไวต่อยาแบบครบถ้วน ในบริบทของห้องปฏิบัติการคลินิกในประเทศไทย มาตรฐานแห่งชาติที่ดูแลโดยกรมวิทยาศาสตร์การแพทย์ สถาบันวิจัยวิทยาศาสตร์สาธารณสุข (NIH Thailand) \cite{dmsc_nih} มีความสอดคล้องอย่างใกล้ชิดกับแนวทางของ CLSI ทำให้ CLSI เป็นข้อมูลอ้างอิงหลักสำหรับสภาพแวดล้อมการติดตั้งเป้าหมาย ฐานข้อมูล Breakpoint ถูกจัดเก็บเป็นตาราง SQLite ที่มีโครงสร้างตาม Schema ดังนี้: ชื่อยาปฏิชีวนะ, กลุ่มเชื้อ, มาตรฐานการทดสอบ (CLSI หรือ EUCAST), กำลังของแผ่นยา (Potency), ขนาดเส้นผ่านศูนย์กลางเกณฑ์ไว ($S_{\min}$ ในหน่วยมิลลิเมตร), ช่วงกึ่งไว ($I_{\min}$--$I_{\max}$ ในหน่วยมิลลิเมตร ถ้ามี) และขนาดเส้นผ่านศูนย์กลางเกณฑ์ดื้อ ($R_{\max}$ ในหน่วยมิลลิเมตร) โครงสร้างนี้เข้ารหัสข้อมูลจากตาราง Breakpoint ที่ตีพิมพ์จากทั้งสองมาตรฐานโดยตรง

\begin{figure}[H]
  \centering
  \safeimg[width=\textwidth]{figures/_review/database_schema.png}{Database ER diagram}
  \caption{แผนภาพ Entity-Relationship (ER) ของฐานข้อมูล SQLite แสดงความสัมพันธ์ระหว่างตาราง
    Users, AnalysisBatch, Plates, PlateResults, Breakpoints\_DiskDiffusion, Microbes,
    Standards และ Antibiotics}
  \label{fig:db_schema}
\end{figure}

สำหรับคู่แผ่นยาและบริเวณยับยั้งแต่ละคู่ที่ตรวจจับได้ ระบบจะดำเนินการตามตรรกะการจำแนกดังนี้:
\begin{enumerate}
  \item ดึงแถว Breakpoint ที่ตรงกับยาปฏิชีวนะที่ระบุได้ด้วย OCR, กลุ่มเชื้อที่ผู้ใช้กำหนด และมาตรฐานที่เลือก (CLSI หรือ EUCAST)
  \item เปรียบเทียบเส้นผ่านศูนย์กลาง ZOI ที่สอบเทียบแล้ว $d_{\mathrm{mm}}$ กับเกณฑ์ที่ดึงมา
  \item กำหนดการจำแนก: ไวต่อยา (S) หาก $d_{\mathrm{mm}} \geq S_{\min}$, ดื้อยา (R) หาก $d_{\mathrm{mm}} \leq R_{\max}$, กึ่งไว (I) หากไม่ตรงกับเงื่อนไขใดข้างต้น (ในกรณีที่มีหมวดหมู่ Intermediate สำหรับคู่ยาและเชื้อนั้น)
  \item หากไม่พบข้อมูล Breakpoint ที่ตรงกันสำหรับคู่ยาปฏิชีวนะและเชื้อที่ระบุ ผลการจำแนกจะถูกส่งคืนเป็น ``ไม่พบข้อมูล (Not Available หรือ NA)'' แต่จะยังคงรายงานเส้นผ่านศูนย์กลางดิบเพื่อให้สามารถค้นหาด้วยตนเองได้
\end{enumerate}

ผลลัพธ์สำหรับคู่แผ่นยาและบริเวณยับยั้งทั้งหมดบนจานจะถูกรวบรวมและส่งคืนไปยัง Frontend ในรูปแบบวัตถุ JSON ที่มีโครงสร้าง ซึ่ง Dashboard ของ React จะนำไปเรนเดอร์เป็นรายงานแบบตารางที่ระบุชื่อยาปฏิชีวนะ, กำลังของแผ่นยา, เส้นผ่านศูนย์กลางบริเวณยับยั้งที่ตรวจพบ (หน่วยมิลลิเมตร) และการจำแนก S/I/R สำหรับแต่ละแผ่นยา นอกจากนี้ Dashboard ยังแสดงภาพซ้อนทับ (Overlay) ของรูปภาพจานเพาะเชื้อต้นฉบับพร้อมกับวาด Bounding box รอบแผ่นยาและบริเวณยับยั้งทั้งหมดที่ตรวจพบ เพื่อให้การยืนยันทางสายตาว่าการตรวจจับมีความถูกต้องเชิงพื้นที่ ผลลัพธ์ทั้งหมดจะถูกบันทึกไว้ในฐานข้อมูลประวัติ SQLite พร้อมกันเพื่อการตรวจสอบและการพิจารณาย้อนหลัง
